\documentclass[11pt,a4paper]{article}
% Common preamble for the hint sheets (\input from each file)
\usepackage[utf8]{inputenc}
\usepackage[T1]{fontenc}
\usepackage{lmodern}
\usepackage{amsmath,amssymb,amsthm,mathtools}
\usepackage[a4paper,margin=2.3cm]{geometry}
\usepackage{enumitem}
\usepackage{xcolor}
\usepackage[hidelinks]{hyperref}
\renewcommand{\leq}{\leqslant}
\renewcommand{\le}{\leqslant}
\renewcommand{\geq}{\geqslant}
\renewcommand{\ge}{\geqslant}
\setlength{\parindent}{0pt}
\setlength{\parskip}{0.35em}
% \hintstitle{Sheet title}{Class work / Homework}
\newcommand{\hintstitle}[2]{\begin{center}{\LARGE\bfseries Hints: #1\par}\vspace{0.5em}{\large #2\par}\end{center}
\noindent\textit{The hints to each problem are cumulative. Read one, try again, and open the next only if you are still stuck.}\par\medskip\hrule\medskip}
\newcommand{\problem}[1]{\par\bigskip{\large\bfseries Problem #1}\par\nopagebreak\smallskip}
\newcommand{\subpart}[1]{\par\smallskip\textbf{(#1)}\par\nopagebreak}
\newcommand{\hint}[1]{\par\smallskip\noindent\textcolor{black!70}{\textbf{Hint #1.}}\enspace}
\newcommand{\note}[1]{\par\smallskip\noindent\textit{Note.} #1\par}
\providecommand{\lt}{<}
\providecommand{\gt}{>}
\providecommand{\tightlist}{\setlength{\itemsep}{0pt}\setlength{\parskip}{0pt}}
\emergencystretch=3em

\begin{document}
\hintstitle{Limits -- 3. Asymptotics of sequences}{Homework}

\problem{1}
\hint{1} The map $g(x)=x-\arcsin(\sin^3x)$ is odd, so it suffices to take $x_0\in(0,1]$ (for $x_0=0$ the sequence is $0$). For $x\in(0,1]$: $0<\sin^3x<\sin x$, so $0<\arcsin(\sin^3x)<x$; hence $x_n$ stays in $(0,1]$, decreases, and tends to a fixed point, i.e.\ to $0$.
\hint{2} Here $\arcsin(\sin^3x)=x^3+O(x^5)$, so $x_{n+1}=x_n-x_n^3+O(x_n^5)$: Theorem 1 with $p=3$, $a=1$. Telescoping: $\dfrac1{x_{n+1}^2}-\dfrac1{x_n^2}\to2$, and Stolz gives $nx_n^2\to\frac12$.
\hint{3} Answer: $\lim\sqrt n\,x_n=\frac1{\sqrt2}$ for $x_0\in(0,1]$, $-\frac1{\sqrt2}$ for $x_0\in[-1,0)$, and $0$ for $x_0=0$.

\problem{2}
\hint{1} Substitute $t=nx$: $b_n=\frac1n\int_0^nf(t)\,dt=\frac1n\sum_{k=0}^{n-1}a_k$.
\hint{2} So $b_n$ is the Ces\`aro mean of $(a_k)$, and converges to the same limit (Stolz--Ces\`aro).

\problem{3}
\hint{1} The recurrence is $\frac{x_{n+1}}{x_n}=2018\cdot\frac{n+1}n$; multiply these ratios.
\hint{2} $x_n=2018^n\,n$, so $n-\log_{2018}x_n=-\log_{2018}n$, and the limit is $\dfrac{\log_2n}{-\log_{2018}n}=-\log_22018$.

\problem{4}
\hint{1} Compute the first terms: $1,3,2,\frac32,\frac74,\frac{13}8,\frac{25}{16},\frac{51}{32},\dots$ and keep track of which of the three previous terms is discarded (the largest).
\hint{2} From $(b_5,b_6,b_7)$ on, the pattern is periodic with period $3$: if $p>q>r$ are three consecutive terms, the next three are $s=\frac{q+r}2$, $t=\frac{q+3r}4$, $u=\frac{q+7r}8$, and again $s>t>u$. So the pair $(q,r)$ is mapped to $(t,u)$, independent of $p$.
\hint{3} The map $(q,r)\mapsto\bigl(\frac{q+3r}4,\frac{q+7r}8\bigr)$ preserves $q+6r$ and divides $q-r$ by $8$. Starting from $(b_6,b_7)=(\frac{13}8,\frac{25}{16})$, $q+6r=11$, so $\lim b_n=\frac{11}7$.

\problem{5}
\hint{1} The candidate is $L=4$ ($L=2\sqrt L$). Note $c_{n+1}$ lies between $2\sqrt{\min(c_n,c_{n-1})}$ and $2\sqrt{\max(c_n,c_{n-1})}$, and $2\sqrt x\ge x$ exactly for $x\le4$.
\hint{2} Hence with $m_n=\min(c_n,c_{n-1},4)$ and $M_n=\max(c_n,c_{n-1},4)$ we get $c_{n+1}\in[m_n,M_n]$: the sequence is bounded and bounded away from $0$.
\hint{3} Let $M=\limsup c_n$, $m=\liminf c_n$ (both in $(0,\infty)$). Passing to the limit, $M\le2\sqrt M$ and $m\ge2\sqrt m$, so $M\le4\le m$; thus $c_n\to4$.

\problem{6}
\hint{1} Rewrite: $x_{n+2}-2x_{n+1}=(x_{n+1}-2x_n)+\frac{2^n}{x_n}$. So $d_n=x_{n+1}-2x_n$ satisfies $d_1=0$, $d_{n+1}=d_n+\frac{2^n}{x_n}$.
\hint{2} Put $y_n=x_n/2^n$. Dividing by $2^{n+1}$: $y_{n+1}-y_n=\frac{d_n}{2^{n+1}}$, and $\frac{2^n}{x_n}=\frac1{y_n}$, so $d_n=\sum_{k=1}^{n-1}\frac1{y_k}$. Hence $d_n\ge0$, $y_n$ is non-decreasing, $y_n\ge y_1=1$ and $d_n\le n-1$.
\hint{3} Then $0\le y_{n+1}-y_n\le\frac{n-1}{2^{n+1}}$ is summable, so $L=\lim y_n$ exists and $L\le1+\sum_{n\ge1}\frac{n-1}{2^{n+1}}=\frac32$. Using $y_k\le\frac32$ in $d_n$ gives $L\ge1+\frac23\cdot\frac12=\frac43$.
\hint{4} Iterating these bounds (compute the first few $y_k$ exactly, estimate the tail) sharpens them; numerically $L\approx1.4782$.

\problem{7}
\hint{1} Factor: $(2n)^4-1=\bigl(1-\frac1{(2n)^2}\bigr)\bigl(1+\frac1{(2n)^2}\bigr)(2n)^4$ and similarly for $(2n+1)^4-1$. Then the product splits into products whose values come from the sine product $\frac{\sin\pi x}{\pi x}=\prod_{n\ge1}\bigl(1-\frac{x^2}{n^2}\bigr)$ (and its versions for $\cos$, $\sinh$, $\cosh$).
\hint{2} Needed: $\prod_{n\ge1}\bigl(1-\frac1{4n^2}\bigr)=\frac2\pi$, $\prod_{n\ge1}\bigl(1+\frac1{4n^2}\bigr)=\frac{2\sinh(\pi/2)}{\pi}$, $\prod_{n\ge1}\bigl(1-\frac1{(2n+1)^2}\bigr)=\frac\pi4$, $\prod_{n\ge1}\bigl(1+\frac1{(2n+1)^2}\bigr)=\frac{\cosh(\pi/2)}2$. The remaining factor is $\frac{(2n)^4(n+1)^2}{(2n+1)^4n^2}=\bigl(1-\frac1{(2n+1)^2}\bigr)^2$.
\hint{3} Combining, the product equals $\dfrac{2}{\pi}\tanh\dfrac\pi2\approx0.5839$.

\problem{8}
\hint{1} Invert: $\dfrac1{a_{n+1}}=\dfrac1{a_n}+1+a_n$.
\hint{2} Hence $\frac1{a_n}\ge\frac1{a_0}+n$, so $a_n\to0$; then $\frac1{a_{n+1}}-\frac1{a_n}\to1$ and Stolz gives $\lim na_n=1$.

\problem{9}
\hint{1} $x_n=P_n(x_1)$, where $P_n$ is a polynomial with non-negative coefficients and no constant term; on $(0,\infty)$ it is strictly increasing. Note $x_{n+1}>x_n\iff x_n>1-\frac1n$.
\hint{2} So the condition is: $x_1>0$ and $1-\frac1n<P_n(x_1)<1$ for all $n$. Let $a_n,b_n>0$ be defined by $P_n(a_n)=1-\frac1n$, $P_n(b_n)=1$. Then $P_{n+1}(a_n)=(1-\frac1n)\cdot1<1-\frac1{n+1}$ and $P_{n+1}(b_n)=1+\frac1n>1$, so $a_n$ increases and $b_n$ decreases; nested intervals give existence.
\hint{3} Uniqueness: $P_n(x)/x$ is non-decreasing, so $P_n(b_n)-P_n(a_n)\ge(b_n-a_n)\frac{P_n(a_n)}{a_n}$, i.e.\ $\frac1n\ge(b_n-a_n)\frac{1-1/n}{a_n}$. With $a_n<1$ this gives $b_n-a_n\le\frac1{n-1}\to0$.

\problem{10}
\hint{1} Fix $\varepsilon>0$ and let $E_N=\{x>0:\ |f(nx)|\le\varepsilon\ \text{for all }n\ge N\}$. Each $E_N$ is closed in $(0,\infty)$ (continuity of $f$), and $\bigcup_NE_N=(0,\infty)$ by hypothesis.
\hint{2} By Baire's theorem some $E_N$ contains an interval $(a,b)$, $0<a<b$. Then $|f|\le\varepsilon$ on $\bigcup_{n\ge N}(na,nb)$.
\hint{3} For $n>\frac a{b-a}$ consecutive intervals $(na,nb)$ and $((n+1)a,(n+1)b)$ overlap, so this union contains a half-line $(M,\infty)$. Hence $|f(x)|\le\varepsilon$ for $x>M$.
\end{document}
